%% APPENDIX C -- the velocity-space line mask and its cost.
%% Owner: appendix-triage.  Carries \label{app:linemask} (cited by
%% 04_method.tex) and \label{sec:linecost}.
\section{The velocity-space line mask, and what it costs}
\label{app:linemask}
% Generated by v361_calc.py -- do not edit by hand.
\begin{table}
\centering
\footnotesize
\setlength{\tabcolsep}{4pt}
\caption{Band-by-band cost of the molecular-line exclusion. Union: unique sky frequency searched in the band. Masked: the part of it inside a $\pm\MaskVWidth$\,km\,s$^{-1}$ stellar-frame tube around a catalogued transition of one of the masked species. Left: what a technosignature search can still occupy. The mask is excluded search space and not a veto (\S\ref{sec:linecost}); the loss is concentrated in Band~6, where CO\,$2\to1$ sits.}
\label{tab:maskband}
\begin{tabular}{@{}crrrr@{}}
\hline
Band & Union (GHz) & Masked (GHz) & Left (GHz) & Lost (\%) \\
\hline
3 & 20.64 & 0.272 & 20.37 & 1.3 \\
4 & 6.94 & 0.049 & 6.89 & 0.7 \\
5 & 6.94 & 0.065 & 6.87 & 0.9 \\
6 & 29.45 & 1.374 & 28.07 & 4.7 \\
7 & 26.74 & 1.589 & 25.15 & 5.9 \\
8 & 6.68 & 0.164 & 6.52 & 2.5 \\
9 & 6.94 & 0.454 & 6.48 & 6.5 \\
10 & 13.80 & 0.844 & 12.96 & 6.1 \\
\hline
All & 118.13 & 4.811 & 113.31 & 4.1 \\
\hline
\end{tabular}
\end{table}


Attribution decides the candidate list, and the mask decides attribution, so
the mask's definition and its price are both given here.

\emph{Definition.} A one-channel coincidence test is too tight for
circumstellar gas: a native channel spans a fraction of a km\,s$^{-1}$ while
real emission is shifted by many channels. The tolerance is set by the
kinematic budget of the gas rather than by the instrument --- Keplerian
motion where these molecules survive ($\lesssim$13\,km\,s$^{-1}$ beyond
10\,au) and linewidths $\lesssim$1\,km\,s$^{-1}$ --- and is
$\pm\MaskHalfKms$\,km\,s$^{-1}$ around each catalogued transition. It is
evaluated topocentrically, in the observed sky frame
(Appendix~\ref{app:conventions}), which is wide enough to absorb the
barycentric and systemic terms it does not apply. The result is empirically
conservative: the largest attributed circumstellar component ($\beta$~Pic,
within \BpicMaxAbsStel\,km\,s$^{-1}$ of systemic) sits at \BpicMaskPct{} per
cent of the half-width.

The mask holds \MaskNTransNew{} transitions of \MaskNSpecNew{} species at
full JPL/CDMS precision \citep{Pickett1998,CDMS2005}. The species criterion
is the one objective rule statable in advance: the transitions ALMA's band
definitions are built around --- the CO isotopologues, HCN, HCO$^{+}$, CS,
SiO, CN, H$_2$CO and [C\,\textsc{i}] --- which dominate the emission of
disc-selected fields.

\emph{The half-width cannot have been tuned to remove a candidate.}
Reclassifying all \NCross{} crossings at $\pm20$ to $\pm100$\,km\,s$^{-1}$
costs \MaskBWTwenty{} to \MaskBWHundred{} per cent of gross bandwidth, and
every crossing that narrowing the mask releases as unattributed is an
identified carbon monoxide line of this survey's own: the two rank-passing
unattributed crossings are the same pair at every half-width in that range.
Widening the mask cannot suppress anything further. The species list is
equally robust: the alternative selection of every Splatalogue transition
over the \UnionIntervals{} searched islands with $E_{\rm u}\le\MaskDbEuK$\,K
and $\log A_{ij}\ge\MaskDbLogA$ --- \MaskDbTrans{} transitions of
\MaskDbCarriers{} species --- masks \MaskDbPctUnion{} per cent of the union
against \MaskUnionPct{} per cent here and changes no disposition, because no
masked crossing exceeds its own control ensemble. ``Coincident with a known
molecular line'' is therefore a statement about a catalogue rather than a
binary property, and no plausible widening or narrowing of either the
half-width or the species list moves a window.

\subsection{The cost of the mask}
\label{sec:linecost}

Masking removes \MaskUnionPct{} per cent of the
unique-frequency union and \MaskGrossPct{} per cent of window-summed gross
bandwidth, CN and CO dominating. It is not uniform, so
Table~\ref{tab:maskband} gives it band by band: it reaches \MaskWorstPct{}
per cent in Band~\MaskWorstBand{} and stays below 1.5 per cent in
Bands~3--5. Nor is it shared equally between the two experiments ---
\MaskPctA{} per cent of the Class~A union against \MaskPctB{} per cent in
Class~B --- so the fine-channel experiment pays \MaskPctRatioAB{} times the
fractional price, its coverage being concentrated in the line-tuned windows a
molecular mask is designed to remove. The two do not sum to the total because
the classes overlap in frequency.

\emph{The masked crossings are not near misses, and the unmasked ones are
not near hits.} Of the \NCross{} windows carrying an on-star crossing,
\NSpatial{} also beat their control ring, \MaskCrossInStageOne{} of those
lying inside the mask and \MaskCrossOutStageOne{} outside it;
\MaskCrossInNoFlag{} more lie inside the mask without beating their ring, of
which \MaskRepN{} are on \MaskRepStar{} alone at \MaskRepVLo{} to
\MaskRepVHi\,km\,s$^{-1}$ from CO\,$2{\to}1$, each already rejected by its
own ring. The remaining \NonExcRingGe{} are ordinary noise crossings whose
own control ring reaches or exceeds the stellar value (ring-maximum median
\NonExcRingMed{} against a crossing median \NonExcHitMed). None of them sits
close to the mask edge: the nearest lies \MaskNearestKms\,km\,s$^{-1}$ from
any catalogued transition of a masked species (\MaskNearestStar,
\MaskNearestMult{} times the adopted half-width), and \MaskNoLineN{} have no
such transition in the window at all.

\emph{Masking is a design selection effect, and it has a price we cannot
buy off.} It separates discovery from classification: a carrier at a line
frequency is as detectable as anywhere else, but it cannot be classified,
because the astrophysical prior at a catalogued transition is overwhelming.
Removing those frequencies from the searched domain is the conservative
choice, at the cost that \emph{we place no constraint at all on artificial
emission lying inside the masked regions, at any signal strength}. A better
design searches every frequency and attaches a \emph{lower prior} to a
line-coincident event instead, since these are exactly the frequencies a
communicator wanting to be found near a natural reference might choose
\citep{CocconiMorrison1959}. The released channels support that second
search, and the released products keep a known-line-coincident list: every
on-source threshold crossing inside the masked intervals, with its offset
from each nearby transition.
